\begin{claim}{C.73}{Proved}{The exact prime-end comparison for arbitrary simply connected frontiers}
Let $O\subset S^2$ be a simply connected domain with
nondegenerate frontier, and fix a conformal map
$\phi:\mathbb D\to O$. Retain the actual groups
\[
 K_O=\{h\in\operatorname{Homeo}^+(S^2):
                    h|_{S^2\setminus O}=\id\},
             \qquad \mathcal D_\partial
                 =\operatorname{Homeo}(\overline{\mathbb D},S^1).
                                                               \tag{C.73.1}
\]
Every $h\in K_O$ induces the identity on the prime-end
circle of $O$. Its actual prime-end extension defines a
continuous injective homomorphism
\[
 \begin{split}
 \iota_O:K_O&\longrightarrow\mathcal D_\partial,\\
 \iota_O(h)|_{\mathbb D}&=\phi^{-1}h\phi .
 \end{split}                                                   \tag{C.73.2}
\]
All groups have the compact-open and inverse topologies.
The image of (C.73.2) is dense. For locally connected
frontier it is the topological isomorphism of C.68.
For every non-locally-connected frontier its image is
proper, and its inverse on that image is not continuous.
The proof supplies the actual sequences and a disk map
witnessing these two failures.

Every actual ambient operator preserving $O$ induces
its own prime-end map, and (C.73.2) intertwines
conjugation by the original operator with conjugation
by that prime-end map. In particular it identifies
an ambient relative centralizer with the intersection
of its image and the corresponding disk centralizer,
with the original operator labels and restriction
kernel retained. This is a continuous injective
comparison; no injection on fundamental groups or
descent of a disk null-homotopy is asserted.
\end{claim}

\begin{proof}
\emph{The prime-end space and the precise access data.}
Let $O^{\rm P}$ denote the prime-end compactification.
The exact extension of $\phi$ is a homeomorphism
$\phi^{\rm P}:\overline{\mathbb D}\to O^{\rm P}$.
The crosscut-chain definition and the neighborhood
basis are those in
\cite[\S3.1]{C_BinderRojasYampolsky2014}.
We use the spherical version, with diameters taken
in the compact sphere. The explicit spherical
definition and functorial action are also recorded
in \cite[\S2, pp.~1571--1572]{C_OrtegaRuizDelPortal2011}.
No continuous landing map on the whole prime-end
circle is assumed.

Here are the exact crosscut facts needed below.
A genuine crosscut is an arc whose interior lies
in $O$ and whose two distinct endpoints lie on
$\partial O$. In the prime-end disk its closure
is an arc with two distinct prime-end endpoints;
their principal sets are the respective singleton
endpoint sets. Each side of the crosscut determines
the corresponding open circle interval.
An access arc landing at an actual frontier point
has a unique endpoint in the prime-end disk, and
that endpoint has the singleton principal set
consisting of the landing point. Equivalently the
radial image under $\phi$ lands at that point.
These precise statements appear in
\cite[p.~1575 and p.~1581, before the proof proper
of Lemma~7]{C_OrtegaRuizDelPortal2011}; they apply to
spherical frontier points, including infinity.

Let $A\subset S^1$ be the set of accessible prime
ends in the specified coordinate. For $a\in A$,
write $\ell(a)$ for its unique principal point,
so
\[
                  \lim_{r\uparrow1}\phi(ra)=\ell(a).
                                                               \tag{C.73.3}
\]
Every nonempty open circle interval $I$ contains
two members of $A$ with distinct values of $\ell$.
To see this without postulating continuity of
$\ell$ on $I$, first choose a prime-end open
neighborhood whose closure meets the circle
only inside $I$. The crosscut-neighborhood basis
supplies a crosscut side contained in this
neighborhood. The two prime-end endpoints of its
closure lie in $I$, and the preceding genuine
crosscut statement gives them distinct actual
landing points. In particular $A$ is dense,
and $\ell|_{A\cap I}$ is never constant.

An ambient homeomorphism $h$ preserving $O$
takes null chains to null chains: it preserves
crosscuts, their sides and cofinal inclusions,
and uniform continuity on $S^2$ preserves
vanishing diameters. Its inverse gives the
inverse action. This is the actual induced
homeomorphism $h^{\rm P}$ of $O^{\rm P}$.
For $h\in K_O$, let $u$ be its action on the
circle through $\phi^{\rm P}$. It preserves
orientation. Indeed its interior action preserves
the oriented relative class of the closed disk;
naturality of the boundary isomorphism to the
oriented circle class gives degree $+1$ on the
boundary. The access statements give
\[
                    u(A)=A,\qquad \ell(ua)=\ell(a).
                                                               \tag{C.73.4}
\]
The second equality uses the actual fact that
$h$ fixes every point of $\partial O$.

\emph{The unlinked access fibres and circle rigidity.}
There cannot be four distinct accessible prime
ends $a,b,c,d$ in cyclic order such that
\[
                   \ell(a)=\ell(c)=p,\qquad
                   \ell(b)=\ell(d)=q,\qquad p\ne q.
                                                               \tag{C.73.5}
\]
For a complete separation argument, take radial
tails at $a$ and $c$ in the closed unit disk
and join their interior endpoints by a simple
arc inside the disk, disjoint from those tails
except at the two joining endpoints. For example,
an arc on a smaller concentric circle gives
such a joining arc. Their union is a genuine
disk crosscut $\alpha$ from $a$ to $c$.
By (C.73.3), $\phi(\alpha\cap\mathbb D)$
closes at both ends to the same point $p$.
It is injective on its interior and does not
contain $p$ there. The resulting map from the
circle obtained by identifying the two endpoints
is a continuous bijection onto its image.
Compactness makes this image a Jordan curve
$J\subset S^2$, and
\[
                         J\cap\partial O=\{p\}.
                                                               \tag{C.73.6}
\]
The two sides of $\mathbb D\setminus\alpha$
map to opposite components of $S^2\setminus J$.
This can be checked at any interior point of
$\alpha$, where the open-set homeomorphism
$\phi$ takes its two local sides to the two
local Jordan sides; connectedness then gives
the same assignment for each entire component.
The radial tails at $b$ and $d$ lie on opposite
disk sides. Their landing point $q$ therefore
lies in the closures of both Jordan domains,
whose intersection is $J$. Equation (C.73.6)
would force $q=p$. This proves the impossibility
of (C.73.5), without any assertion about other
impressions.

Suppose $u$ is nonidentity. Its moved set is
open and meets $A$, so choose a moved $a\in A$.
Choose real representatives so that
$a<u(a)<a+1$. There is a small open interval
$I$ just after $a$ such that every $b\in I$
satisfies
$a<b<u(a)<u(b)<a+1$.
The nonconstancy proved above supplies
$b\in A\cap I$ with $\ell(b)\ne\ell(a)$.
Equation (C.73.4) now gives exactly the forbidden
four-tuple $a,b,u(a),u(b)$.
Consequently $u=\id$.
This proves the prime-end rigidity assertion
and defines (C.73.2) on every element of $K_O$.
The group equations follow from the interior
conjugacy and density of the interior. Injectivity
also follows from that conjugacy and the specified
identity outside $O$.

\emph{Continuity in the actual ambient topology.}
The group $K_O$ is a closed subgroup of
$\operatorname{Homeo}^+(S^2)$, hence Polish;
$\mathcal D_\partial$ is Polish as well.
The graph
\[
       \mathcal G=
       \{(h,g)\in K_O\times\mathcal D_\partial:
                       g|_{\mathbb D}=\phi^{-1}h\phi\}
                                                               \tag{C.73.7}
\]
is closed. Namely, convergence of both coordinates
preserves the displayed equality at every interior
point; density and the fixed circle then determine
$g$ uniquely. Thus $\mathcal G$ is a closed
Polish subgroup. Its projection to $K_O$ is a
continuous bijective homomorphism, by the existence
and uniqueness already proved. The Polish-group
open-mapping consequence of the orbit theorem
used in C.30 makes that projection open: apply
the theorem to the transitive left action of
$\mathcal G$ on the Polish group $K_O$ through
the projection. Its inverse is therefore
continuous. Composing that inverse with the
second projection proves continuity of $\iota_O$,
including for compact families of maps and
inverses. This proves no continuity of its
inverse on its image.

Every compactly supported disk homeomorphism
$g$ belongs to the image. Its interior conjugate
$\phi g\phi^{-1}$ has compact support in $O$,
so extension by the identity gives an element
of $K_O$. Such disk maps are dense in
$\mathcal D_\partial$: for $0<r<1$, define
\[
 g^{[r]}(z)=
 \begin{cases}
      r\,g(z/r),& |z|\leq r,\\
      z,& |z|\geq r .
 \end{cases}                                                   \tag{C.73.8}
\]
The boundary-fixed equation makes these formulas
agree. Their inverse has the same formula with
$g^{-1}$. Uniform continuity of $g,g^{-1}$
on the closed disk gives
$g^{[r]}\to g$ with inverses as $r\uparrow1$.
Thus the image in (C.73.2) is dense.

\emph{A nontrivial impression gives both failures of comparison.}
In the fixed coordinate the impression of
$\zeta\in S^1$ can be written
\[
 I(\zeta)=\bigcap_{\epsilon>0}
       \overline{\phi\bigl(\mathbb D\cap
                                  B(\zeta,\epsilon)\bigr)}
                       \quad\text{in }S^2 .
                                                               \tag{C.73.9}
\]
This is the usual impression: the disk
neighborhoods and the crosscut neighborhoods
are cofinal bases in $O^{\rm P}$.
The sets in (C.73.9) are nested compact
connected sets, and their intersection is
nonempty and contained in $\partial O$.

A non-locally-connected frontier has a
nontrivial impression. Here is also a direct
verification of this consequence. If every
$I(\zeta)$ were a singleton, assign that value
to $\phi(\zeta)$. Any interior sequence tending
to $\zeta$ would have all its spherical cluster
points in that singleton, so would converge
to it. For a sequence of boundary points
$\zeta_n\to\zeta$, approximate each assigned
value by an interior value at a point within
$1/n$ of $\zeta_n$. The same interior-sequence
argument then proves continuity on the circle
and on the whole closed disk. Every frontier
point is a value on the circle, by compactness
applied to an interior approximating sequence.
The frontier would therefore be a continuous
image of the compact locally connected circle,
hence locally connected, by
\cite[Lemma~3.2]{C_BinderRojasYampolsky2014}.
This contradiction proves the assertion, and
also keeps the spherical topology explicit.

Fix distinct $a,b\in I(\zeta)$.
Choose successively positive numbers
$\epsilon_n\to0$ and points
$u_n,v_n\in\mathbb D\cap B(\zeta,\epsilon_n)$
such that
\[
                    \phi(u_n)\longrightarrow a,\qquad
                    \phi(v_n)\longrightarrow b .
                                                               \tag{C.73.10}
\]
Their existence follows from (C.73.9).
The lens $\mathbb D\cap B(\zeta,\epsilon_n)$
is convex and open, so a small closed
neighborhood of the segment $[u_n,v_n]$
is a convex topological disk $D_n$ contained
in that lens, with both points in its interior.
After choosing $D_n$, choose the next
$\epsilon_{n+1}$ smaller than its distance
from $\zeta$ and the distances of all earlier
disks from $\zeta$. This retains (C.73.10)
and makes the $D_n$ pairwise disjoint.
They form a locally finite family in
$\mathbb D$, with sole possible accumulation
point $\zeta$ in the closed disk, and
$\operatorname{diam}D_n\leq2\epsilon_n$.

Choose a homeomorphism $g_n$ supported in
$D_n$ taking $u_n$ to $v_n$. There is a
concrete such map: on each segment from
$u_n$ to a boundary point $z\in\partial D_n$,
send $u_n+t(z-u_n)$ to
$v_n+t(z-v_n)$, $0\leq t\leq1$.
Convexity and interiority give the unique
ray coordinates and the inverse formula
with $u_n,v_n$ interchanged. It fixes
$\partial D_n$ and extends by the identity
outside $D_n$. Therefore
\[
 \sup|g_n-\id|\leq2\epsilon_n,\qquad
 \sup|g_n^{-1}-\id|\leq2\epsilon_n .
                                                               \tag{C.73.11}
\]
Let $h_n\in K_O$ be its compactly supported
ambient conjugate. Then $\iota_O(h_n)=g_n$
tends to the identity in $\mathcal D_\partial$,
but $h_n$ does not tend to the identity:
it moves $\phi(u_n)$ to $\phi(v_n)$, whose
limiting points in (C.73.10) are distinct.
This proves the claimed failure of inverse
continuity on the actual image.

There is also one map witnessing nonsurjectivity.
Set $g=g_n$ on each $D_n$ and set $g=\id$
elsewhere on $\overline{\mathbb D}$.
Local finiteness gives a homeomorphism in the
interior, and the shrinking supports give
continuity at $\zeta$ and every other boundary
point. The same reasoning applies to the
piecewise inverses, so $g\in\mathcal D_\partial$.
If $g=\iota_O(h)$ for an $h\in K_O$, the
interior conjugacy would give
$h(\phi(u_n))=\phi(v_n)$.
Continuity at the fixed frontier point $a$
would then imply $a=b$, a contradiction.
Thus this specified $g$ is outside the image.
Together with density this proves the stated
proper-image assertion. For locally connected
frontier, C.68 gives the full opposite
comparison and its topology; no such
comparison is inferred here at a nontrivial
impression.

\emph{Original operators and the remaining filling problem.}
Let $H_O$ be any actual ambient operator group
preserving $O$. Denote its same-label
prime-end action in the chosen coordinate by
\[
       \widehat h=(\phi^{\rm P})^{-1}
                           h^{\rm P}\phi^{\rm P}
               \quad(h\in H_O).
                                                               \tag{C.73.12}
\]
An operator may reverse orientation. Its
conjugation still preserves $K_O$, and the
interior equations and density give exactly
\[
       \iota_O(hkh^{-1})
           =\widehat h\,\iota_O(k)\,\widehat h^{-1}
                       \quad(k\in K_O,\ h\in H_O).
                                                               \tag{C.73.13}
\]
Consequently
\[
 \iota_O\bigl(C_{K_O}(H_O)\bigr)
       =\iota_O(K_O)\cap
                 C_{\mathcal D_\partial}(\widehat H_O).
                                                               \tag{C.73.14}
\]
For the reverse implication, the disk equation
first gives the ambient commutation on $O$;
on its complement both relative maps are
the identity. This proves the original
same-label equation everywhere. An operator
acting trivially on $O$ induces the identity
prime-end map and already commutes with every
element of $K_O$, so the restriction kernel
is retained by (C.73.13).

The density proved above concerns
$\iota_O(K_O)$, not an unproved density
statement for the centralizer intersection.
Likewise a disk contraction need not stay
in this image, and (C.73.11) shows why
continuity of a putative inverse transfer
requires its own argument. These exact
comparisons leave the general equivariant
frontier-filling obligation of C.12 open.
\end{proof}
